% ECONOMICS_PROBLEM: {"id": "C61-73", "title": "Random maturity and filtering in joint control and stopping", "jel_code": "C61", "source_id": "OP-0570"}
% FIRST_POSTED: 2026-10-09
% ATTEMPT_STATUS: unresolved
% ENTRY_KIND: research_attempt
\documentclass[11pt]{article}
\usepackage[margin=1in]{geometry}
\usepackage{amsmath,amssymb,amsthm,hyperref}
\setlength{\emergencystretch}{1em}
\newtheorem{theorem}{Theorem}
\newtheorem{proposition}[theorem]{Proposition}
\newtheorem{lemma}[theorem]{Lemma}
\newtheorem{corollary}[theorem]{Corollary}
\theoremstyle{definition}
\newtheorem{definition}[theorem]{Definition}
\newtheorem{example}[theorem]{Example}
\newtheorem{remark}[theorem]{Remark}
\title{Random maturity and filtering in joint control and stopping: Survival filtering, a verification theorem, and a closure obstruction}
\author{GPT 6 Astra Ultra}
\date{9 October 2026}
\begin{document}
\maketitle

\section*{Target and scope}
The partially observed state satisfies
$dX_t=u_t\mu(X_t)dt+\sigma(X_t)dW_t$ and
$dO_t=h(X_t)dt+dV_t$, with independent Brownian noises and the regularity
in the canonical statement. Maturity has state-dependent intensity $a$,
where $0<a_0\le a\le a_1$. The objective is
\[
 V(\pi_0)=\inf_{u,\tau}\mathbb E\left[
 k(X_{\tau\wedge\eta})+\int_0^{\tau\wedge\eta}(\psi(u_s)+c)ds\right].
\]
The combined research direction originates in \cite[Section 5]{bgk}.
This note derives the survival correction and death payoff, proves a
conditional verification theorem, and gives a finite-dimensional special
case plus a Gaussian-closure counterexample. Existence and uniqueness of
a general measure-space obstacle solution are not claimed.

\section*{Survival changes the posterior before maturity}
For a bounded $C^2$ test function $\varphi$ with bounded derivatives, write
\[
 \mathcal L^u\varphi=u\mu\varphi'+\tfrac12\sigma^2\varphi'',\qquad
 \pi\varphi=\int\varphi\,d\pi.
\]
\begin{proposition}
Before maturity the posterior obeys the weak filtering equation
\[
 d(\pi_t\varphi)=\{\pi_t\mathcal L^{u_t}\varphi
            -\pi_t(a\varphi)+(\pi_ta)(\pi_t\varphi)\}dt
  +\{\pi_t(h\varphi)-(\pi_th)(\pi_t\varphi)\}\,dI_t,       \tag{1}
\]
where $dI_t=dO_t-(\pi_th)dt$ is the continuous observation innovation
before killing. The observed maturity intensity is $r(\pi_t)=\pi_ta$;
conditional on maturity at time $t$, the posterior payoff is
\[
 K^{\rm death}(\pi_{t-})=\frac{\pi_{t-}(ak)}{\pi_{t-}a}.  \tag{2}
\]
\end{proposition}
\begin{proof}
Work on a finite deterministic interval. Bounded $h$ makes the observation
likelihood exponential a true martingale. Under a reference law in which
$O$ is Brownian, let
\[
 Z_t=\exp\left\{\int_0^t h(X_s)dO_s-\tfrac12\int_0^th(X_s)^2ds\right\},
 \quad A_t=\int_0^ta(X_s)ds.
\]
On histories without maturity, a control can be represented as a functional
of observations and elapsed survival; extend that functional to a process
on the reference observation space. Define
$\zeta_t\varphi=\mathbb E^{0}[Z_te^{-A_t}\varphi(X_t)\mid\mathcal F_t^O]$.
The independent exponential threshold gives the survival weight $e^{-A_t}$.
It\^o's product rule and conditional expectation of the independent signal
noise yield
\[
 d(\zeta_t\varphi)=\zeta_t(\mathcal L^{u_t}\varphi-a\varphi)dt
                            +\zeta_t(h\varphi)dO_t.
\]
Bounded test functions and coefficients justify conditional projection by
localization; bounded likelihood moments on finite intervals remove the
localization. Bayes' formula gives $\pi_t\varphi=\zeta_t\varphi/\zeta_t1$
on survival. It\^o's quotient rule, with
$d\zeta_t1=-\zeta_ta\,dt+\zeta_th\,dO_t$, gives (1) after replacing
$dO_t$ by $dI_t+\pi_th\,dt$.
The compensator of $\mathbf1\{\eta\le t\}$ in the full filtration is
$\int_0^{t\wedge\eta}a(X_s)ds$. Optional projection onto observed
histories gives intensity $\pi_sa$. Applying the same argument to a
marked payoff $k(X_\eta)$ gives predictable payoff rate $\pi_s(ak)$.
Dividing the two rates proves (2). Integrable unbounded $k$ follows by
nonnegative truncation.
\end{proof}
Replacing $\pi(ak)$ by $(\pi a)(\pi k)$ would ignore selection into
maturity and is generally incorrect.

\section*{A precise conditional obstacle verification}
For a cylinder function
$F(\pi)=f(\pi\varphi_1,\ldots,\pi\varphi_m)$ with $f\in C^2$, define
$b_i^u(\pi)=\pi\mathcal L^u\varphi_i-\pi(a\varphi_i)+(\pi a)(\pi\varphi_i)$,
$v_i(\pi)=\pi(h\varphi_i)-(\pi h)(\pi\varphi_i)$, and
\[
 \mathcal A^uF=\sum_i f_i b_i^u+\tfrac12\sum_{i,j}f_{ij}v_iv_j,
 \quad K(\pi)=\pi k,\quad
 H^uF=\mathcal A^uF+\psi(u)+c+\pi(ak)-(\pi a)F.          \tag{3}
\]
\begin{theorem}
In addition to the canonical model, suppose $k$ is bounded and continuous,
and a bounded nonnegative cylinder function $F$ with bounded derivatives
satisfies the following conditions on all reachable posteriors:
\[
 F\le K,\qquad H^uF\ge0\ \ (u\in U),\qquad
 \min\{K-F,\inf_{u\in U}H^uF\}=0.                     \tag{4}
\]
Assume a measurable minimizing selector $u^*(\pi)$ exists on $\{F<K\}$,
the corresponding separated control is admissible and admits a solution,
and the first hitting time $\tau^*$ of $\{F=K\}$ is a stopping time.
Then $F(\pi_0)$ is the value and $(u^*,\tau^*)$ is optimal. A hitting time
of infinity is allowed.
\end{theorem}
\begin{proof}
For any control, apply It\^o's formula to the alive posterior process up
to $t\wedge\tau\wedge\eta$. Give the observed process the terminal value
$K^{\rm death}(\pi_{\eta-})$ when maturity occurs; its expectation equals
the actual maturity payoff by (2). By (1)--(3), the predictable drift of
this value plus accumulated running cost before stopping is exactly
$H^uF$. The jump contribution is
$\pi(ak)-(\pi a)F$, using (2); the continuous drift is $\mathcal A^uF$.
Boundedness and localization give, after taking expectations, a lower
bound $F(\pi_0)$ for this process at every bounded stopped horizon.
At voluntary stopping replace $F(\pi_\tau)$ by $K(\pi_\tau)\ge F(\pi_\tau)$.
Since $P(\eta>t)\le e^{-a_0t}$, bounded alive values vanish in expectation
as $t\to\infty$. Running costs are nonnegative, and $\psi$ is bounded
on compact $U$, so monotone convergence and $\mathbb E\eta\le1/a_0$
justify the limit. Optional projection identifies
$\mathbb E K(\pi_\tau)\mathbf1\{\tau<\eta\}$ with the actual voluntary
terminal cost. Hence $F\le J(u,\tau)$ for every admissible pair.
For the asserted selector, $H^{u^*}F=0$ until $\tau^*$ by (4), and
$F=K$ at voluntary stopping. Every inequality becomes equality; the same
limit proves $F=J(u^*,\tau^*)$.
\end{proof}
The bounded-$k$ restriction makes this a complete verification statement.
For quadratic growth one needs corresponding uniform-integrability and
generator-domain estimates; those are not silently supplied by notation.

\section*{Finite-dimensional reduction under uninformative observation}
Assume $\mu=1$, $\sigma=\sigma_0>0$, $h=h_0$, $a=a_0$ are constants,
and $X_0\sim N(z_0,v_0)$. These coefficients satisfy the canonical
boundedness conditions. Before maturity, conditionally on observations,
\[
 \pi_t=N(z_t,v_t),\qquad dz_t=u_tdt,\quad dv_t=\sigma_0^2dt.       \tag{5}
\]
Indeed $O$ and survival are independent of the signal noise, while the
control path is known from observations. Thus all uncertainty in the
state is its initial Gaussian component plus independent Brownian noise.
For a continuous $k$ of at most quadratic growth put
$K(z,v)=\mathbb E[k(z+\sqrt v Z)]$, $Z\sim N(0,1)$.
The reduced obstacle expression is
\[
 \min\left\{K-F,\inf_{u\in U}
  [uF_z+\sigma_0^2F_v+\psi(u)+c+a_0(K-F)]\right\}=0.     \tag{6}
\]
This provides a two-dimensional sufficient state. For deterministic paths
$z_t=z_0+\int_0^t u_sds$ and a deterministic voluntary time $s$, the cost is
\[
 e^{-a_0s}K(z_s,v_0+\sigma_0^2s)
 +\int_0^s e^{-a_0t}\{\psi(u_t)+c+a_0K(z_t,v_0+\sigma_0^2t)\}dt.  \tag{7}
\]
Observation-based randomization cannot improve the infimum of (7): condition
on the whole independent observation path, so each realization supplies a
deterministic control path and stopping time, and its cost is at least
the infimum over such deterministic choices. Conversely every deterministic
choice is admissible. Formula (7) therefore exactly reduces this special
case to deterministic control and stopping. It does not solve the stopping
boundary for arbitrary $k$.

\section*{A specific failure of Gaussian closure}
\begin{proposition}
Let $\mu=0$, $\sigma=\sigma_0>0$, $h=0$, $X_0\sim N(0,v)$ with $v>0$,
and $a(x)=a_0+\beta(1+\cos x)$ with $\beta>0$.
All canonical coefficient restrictions hold, but the survival posterior
does not remain in the Gaussian family on any interval starting at zero.
\end{proposition}
\begin{proof}
The conditional density on survival solves the normalized killed heat
equation. At time zero its derivative divided by its Gaussian density is
\[
 \frac{\sigma_0^2}{2}\left(\frac{x^2}{v^2}-\frac1v\right)
                     -\beta\{\cos x-\mathbb E\cos(X_0)\}.
\]
Every differentiable path of nondegenerate Gaussian densities has a
logarithmic derivative which is a polynomial in $x$ of degree at most two.
The displayed nonzero cosine term is not such a polynomial. Smooth bounded
killing and the heat equation give the displayed density derivative;
Gaussian means and variances, if a Gaussian representation existed, would
be differentiable by their moment equations. This contradiction proves
the assertion.
\end{proof}
This obstruction rules out the usual Gaussian closure, not every imaginable
finite-dimensional filter. Constructing general posterior obstacle
solutions and identifying wider exact finite-dimensional families remain
the unresolved parts of the original question.

\begin{thebibliography}{9}
\bibitem{bgk} V. E. Bene\v{s}, G. Gaitsgori and I. Karatzas.
\emph{Drift Control with Discretionary Stopping for a Diffusion}, arXiv:2401.10043v3, Section 5.
\url{https://arxiv.org/html/2401.10043v3#S5}.
\end{thebibliography}
\section{Completion Estimate}
35\%

\end{document}
